\scanpage{102}
\mspage{5\deleted{\unsure{1|0}}\unsure{0|1}'}

$\exists!\,\lambda_{e'}\in\Lambda$ tel qu'on ait

\[
\lambda_{\Gamma,s}=\lambda_{e'}\underset{\mu_2}{\wedge}(e\wedge e')\tag*{(71)}
\]

et on a bien sur

\[
\lambda_{\deleted{\unsure{e}}\overline{e'}\deleted{\unsure{e}}}=-\lambda_{e'}\qquad\text{si }\mu(i')=\{e',\overline{e'}\}\tag*{(\unsure{72|71})}
\]

Ainsi, $\{\lambda_{e'},\lambda_{\overline{e'}}\}$ est l'image inverse de $\dot\lambda\in\Lambda^*/\pm1$ dans $\Lambda^*$, et il existe un unique $e'\in\mu(i')$, tel qu'on ait

\[
\lambda=\lambda_{e'}\quad\text{i.e.}\quad\lambda_{\Gamma,s}=\lambda\wedge(e\wedge e')\in\Gamma_{ii'}.\tag*{(73)}
\]

A priori, ce choix de $e'$ dépend de celui de $s\in S$ -- en fait il est indépendant de ce choix. On trouve ainsi l'opération de $\Lambda=\mathbb Z/n\mathbb Z$ sur $\deleted{\mathbb Z}\vec\Pi$, qui fait de $\unsure{\vec\Pi|\overline\Pi|\vec\pi}$ un torseur sous $\Lambda$.

L'application

\[
\tau_F:\underbrace{\vec F}_{\text{\unsure{ou|en} }\vec F_M}\quad\unsure{\flat|\ell}\longrightarrow\vec\Pi\tag*{(74)}
\]

qui relève $\tau:F\longrightarrow\Pi$ est maintenant évidente, car \deleted{\unsure{face}}\deleted{toute} face \inserted{orientée} $f\in\vec F$ définit à la fois un $\tau(f)=i\in\Pi$, et un $e\in\deleted{\unsure{P}}\mu(i)=\operatorname{Or}P_i$ orientation de la structure triangulaire à laquelle appartient cette face.
% archive sheet number: 101
